\documentclass{report}
\usepackage{amsmath,amssymb}
\setlength{\parindent}{0mm}
\setlength{\parskip}{1em}
\begin{document}
\begin{center}
\rule{6in}{1pt} \
{\large Sarah Osborn \\
{\bf A Multilevel Solution Strategy for the Stochastic Galerkin Method for PDEs with Random Input Data}}

Center for Applied Scientific Computing \\ Lawrence Livermore National Laboratory \\ Box 808 \\ L-561 \\ Livermore \\ CA 94551
\\
{\tt osborn9@llnl.gov}\\
Eric Phipps\\
Victoria Howle\end{center}

We discuss solving partial differential equations (PDEs) with random
input data using the stochastic Galerkin method. This method can often be
computationally demanding as the method suffers from the \textit{curse of
dimensionality} where the computational effort increases greatly as the
stochastic dimension increases. We consider a multilevel solution
strategy for the stochastic Galerkin method that employs hierarchies of
spatial and stochastic approximations in an attempt to diminish some of
the computational burden. Analysis of the proposed multilevel method and
numerical results are presented that compare the multilevel approach to
the traditional, single-level stochastic Galerkin method.


\end{document}
